\documentclass{jsarticle}
\usepackage{amsmath}
\usepackage{amssymb}
\begin{document}
\title{M theory equal with AdS5 manifold, \\
Gamma function escort into Beta function}
\author{}
\date{}
\maketitle

   フェルマー型のカラビ・ヤウ多様体が、ゼータ関数を部品にしている。
   $$ x^a + y^b + z^c + u^d + v^e = 0, ({1 \over a} + {1 \over b}
   + {1 \over c} + {1 \over d} + {1 \over e} = 1) $$
   $$ x^a + y^b + z^c + u^d + v^e - 5 \psi xyzuv = 0 $$
   そのオイラー数が、ホッジ数を経て、
   $$ e = 2(h^{ij} - h^{ji}) $$
   大域的積分多様体のガンマ関数となり、　
   $$ \int \Gamma(\gamma)^{'}dx_m = \int \Gamma dx_m \cdot {d \over d \gamma}
   \Gamma \le \int \Gamma dx_m + {d \over d \gamma}\Gamma $$
   Jones多項式を形成して、
   $$ = e^{x /log x} + e^{-x \log x} \ge e^{x \log x} - e^{-x \log x} $$
   鏡映理論となり、
   $$ R_{ij}^{'} = - R_{ij} $$
   ヒッグス場から、
   $$ {d \over df}F = {d \over df}\int\int{{1 \over {(x \log x)^2}}}dx_m 
   + {d \over df}\int\int{{1 \over {(y \log y)^{1 \over 2}}}}dy_m $$
   オイラーの定数の多様体積分を加群として、
   $$ {d \over df}F + \int C dx_m = \int {(\int {1 \over {x^s}}dx 
   - \log x)}\mathrm{dvol} = e^{-x \log x} + e^{x \log x} $$
   リッチテンソルを時間における流体理論として、
   $$ {d \over dt}g_{ij}(t) = - 2 R_{ij} $$
   周期関数となり、
   $$ {d \over df}F + \int C dx_m = 2 (\cos (i x \log x) - i\sin(i x \log x)) $$
   全ては、ホッジ予想となる。
   $$ = 2(h^{ij} - h^{ji}) $$
   全ては、カラビ・ヤウ多様体が、ゼータ関数を部品とする、オイラーの定数の
   多様体積分として、ガンマ関数における大域的積分多様体と同型となり、
   ホッジ数が、５次元型フェルマー方程式における、リーマン予想を基点にする
   D-braneを解にもっていく、共形場理論の鏡映理論となる、ヒッグス場
   方程式が、世界をプラトニックな空間を経て、スピリチュアルな空間と
   物質な空間における、情報の変換を成している、心の影を形成している。
   心理学から物理学、数学へと、情報が行き来している様が、上の式たちである。

   その結果の式が、$AdS_5$多様体である。
   $$ ||ds^2|| = e^{-2\pi T ||\psi||}[\eta_{\mu\nu} + \bar{h}(x)]
   dx^{\mu}dx^{\nu} + T^2 d^2 \psi $$
   それが、
   $\beta(p,q) = \int e^{\sin \theta \cos \theta} \int \sin \theta \cos \theta
   d \theta $
   $$ = {d \over df}F + \int C dx_m$$
   $$ = 2 (\cos (i x \log x) - i \sin(i x \log x)) $$
   と、周期関数へと結論が下る。
   それゆえに、ホッジ予想が解決される。

\end{document}
